\begin{table}[h]
\centering
\caption{Physical role of each data source in peatland fire prediction.}
\label{tab:dataset_role}
\begin{tabular}{p{2.4cm} p{2.9cm} p{5.7cm}}
\toprule
\textbf{Dataset} & \textbf{Predictive role} & \textbf{Fire-relevant physical mechanism} \\
\midrule
FIRMS VIIRS 375~m
    & Label (target)
    & Thermal anomaly at 375~m resolution;
      flaming and near-surface combustion.
      Does \textit{not} detect subsurface smoldering. \\
\addlinespace
ERA5-Land
    & Weather + soil moisture
    & Temperature, humidity, and wind speed control evaporation and fire spread.
      swvl1 (0--7~cm) and swvl2 (7--28~cm) provide modeled surface and shallow peat moisture
      \textit{continuously} at daily resolution. \\
\addlinespace
CHIRPS v3 SAT
    & Rainfall
    & Cumulative rainfall deficit over 2--4 weeks dries the peat surface.
      In equatorial Riau, temperature is near-constant year-round;
      rainfall is the dominant seasonal driver. \\
\addlinespace
Sentinel-1 SAR
    & Measured surface moisture + biomass
    & C-band radar backscatter (VV/VH) penetrates clouds and smoke.
      Wet soil returns stronger signal; dry peat returns weaker signal.
      Directly \textit{measures} moisture, unlike ERA5 which \textit{models} it.
      Sparse cadence (6--12~day) is the trade-off. \\
\addlinespace
Dynamic World
    & Land cover probabilities
    & 9 continuous class probabilities (trees, crops, bare, built, etc.)
      from Sentinel-2 at 10~m.
      Distinguishes fire-prone land covers (plantation, shrubland)
      from non-flammable surfaces (water, built-up).
      Cloud cover causes gaps during fire season. \\
\addlinespace
Peat depth\hfill (BRGM FEG 2019)
    & Static vulnerability
    & Peat thickness determines fire persistence:
      deep peat ($>\!3$~m) can smolder for weeks underground and reignite surface fires.
      Non-peat cells lack this mechanism entirely. \\
\bottomrule
\end{tabular}
\end{table}
